\section{Post-Audit Exact Refinements}\label{sec:v5-refinements}

This section integrates the September 2026 audit into the manuscript rather
than treating it as evidence for a universal conclusion.  The corrected
results sharpen the finite certificate search, add reusable exact-mass
combinators, and isolate the remaining universal obligation.

\subsection{Corrected arithmetic interface}

Let $p\equiv1\pmod4$ be prime and let
\[
  0<c\le2p,\qquad c\equiv3\pmod4,\qquad
  a=\frac{p+c}{4},\qquad M=ap.
\]
Then $a<p$ and $\gcd(c,M)=1$.  Call $p$ \emph{$c$-representable} when
there is a positive divisor $e\mid M^2$ such that
\[
  c\mid M+e
  \quad\text{and}\quad
  c\mid M+M^2/e.
\]
The reconstruction
\[
  x=a,\qquad y=\frac{M+e}{c},\qquad
  z=\frac{M+M^2/e}{c}
\]
is integral and satisfies $1/x+1/y+1/z=4/p$.  Conversely, ordering any
solution $x\le y\le z$, taking $c=4x-p$, and completing the rectangle gives
\[
  (cy-M)(cz-M)=M^2,
\]
so $e=cy-M$ supplies the divisor certificate.  Thus the shift criterion is
an exact reformulation, not merely a sufficient search heuristic.

If $M=\prod_i q_i^{k_i}$, define the signed exponent box
\[
  B(M,c)=\left\{\prod_i q_i^{g_i}\bmod c:-k_i\le g_i\le k_i\right\}.
\]
Because $M$ is a unit modulo $c$, the shift is representable exactly when
$-1\in B(M,c)$.  This statement uses the factorization of $M$, not only that
of $a$; conflating those boxes loses a target congruence.

\subsection{A smaller complete shift range}

For $p>5$, every solution has a least denominator at most $(p-1)/2$.
Indeed, reducing $4xyz=p(xy+xz+yz)$ modulo $p$ shows that some denominator
is divisible by $p$.  If two are $pb$ and $pe$, the remaining denominator
$x$ obeys $4=p/x+1/b+1/e\le p/x+2$, so $x<p/2$.  If exactly one is $pt$,
then $p\mid4t-1$; writing $4t-1=kp$ gives $k\equiv3\pmod4$ and hence
$t\ge(3p+1)/4$.  Were the other two denominators both larger than $p/2$,
their reciprocal sum, including $1/(pt)$, would be at most
\[
  \frac4{p+1}+\frac4{p(3p+1)}<\frac4p,
\]
a contradiction.

It follows that the least-denominator shift satisfies $c=4x-p\le p-2$.
The endpoint $c=p-2$ is impossible for $p>5$ (a proof is given in
\cref{app:v5-proofs}), leaving the complete range
\[
  c\le p-6.
\]
This removes shifts having provably zero ordered accepted mass.  For the
specified sampler that chooses uniformly among $c=3,7,\ldots,2p$ and then
uniformly among divisors of $M^2$, truncation to $c\le p-6$ multiplies its
ordered accepted mass by
\[
  \frac{2(p-1)}{p-5}.
\]
The identity also holds when both masses are zero: it improves a sampler but
does not prove that an accepting divisor exists.

\subsection{Conditional character mass and recursive transport}

For a state $(n,x)$ set $h=4x-n$ and factor $x=\prod_i r_i^{e_i}$.
Choose independent signed exponents $-e_i\le z_i\le e_i$ and write
$R=\prod_i r_i^{z_i}$.  With $\chi$ the Jacobi character modulo $h$, every
Type-II target $R\equiv-1\pmod h$ lies in the event $\chi(R)=-1$.  If this
conditioning event is nonempty, its probability is
\[
  Z=\frac{1-\eta}{2},\qquad
  \eta=\prod_{\chi(r_i)=-1}\frac{(-1)^{e_i}}{2e_i+1},
\]
and $2/5\le Z\le2/3$.  Conditioning therefore preserves all accepted seeds
and multiplies exact success mass by $1/Z$, between $3/2$ and $5/2$.
When no negative-character factor exists, the event is null and the kernel
rejects rather than conditioning on probability zero.

This gain composes with the exact decreasing-state transport.  If $L\ge2$
divides $x$, $y=x/L$, and $q=4y-h\ge2$, then a witness
$(y,qb,qe)$ maps to $(Ly,Lnb,Lne)$.  Both $y<x$ and $q<n/2$, so induction is
well founded.  Using identical branch weights before and after conditioning
gives
\[
  \frac32 U(n,x)\le C(n,x)\le\frac52 U(n,x).
\]
In particular, $C(n,x)>0$ exactly when $U(n,x)>0$: conditioning amplifies
existing mass and cannot manufacture the missing universal positivity
premise.

\subsection{Fail-closed reachability and case unions}

Recursive SWC certificates may include failed terminals and cycles.  A
finite-state certificate supplies $L(s)\in[0,1]$, a nonnegative potential
$V(s)$, and $\varepsilon>0$ such that at nonterminal states
\[
  L(s)\le\mathbb E[L(S_{t+1})\mid S_t=s],\qquad
  V(s)\ge\varepsilon+\mathbb E[V(S_{t+1})\mid S_t=s].
\]
With $L=1$ on success, $L=0$ on failure, and $V=0$ on both terminal types,
\[
  \Pr(\text{success by }H)
  \ge L(s_0)-\frac{V(s_0)}{\varepsilon H}.
\]
Every successful terminal must still carry a checked witness for the original
target.  The lower bound is useful only when its right-hand side is positive.

The accompanying exact-arithmetic kernel checks finite masses, same-target
mixtures, integral witness transport, one- and two-parameter polynomial
families, instantiation, and finite unions of proved case domains.  The
current case certificate contains twelve infinite families and two direct
instances.  Induction over a finite case list proves the union of those cases;
it does not prove that the union contains every prime.

\subsection{Final literature audit and remaining obligation}

The final audit compared these constructions with type-I/type-II average
counts by Elsholtz and Tao
(\href{https://arxiv.org/abs/1107.1010}{arXiv:1107.1010}), the exact divisor
and sieve criteria of Dahan
(\href{https://arxiv.org/abs/2608.24035}{arXiv:2608.24035}), the divisor
parametrization of Bello-Hern\'andez, Benito, and Fern\'andez
(\href{https://arxiv.org/abs/2606.10922}{arXiv:2606.10922}), and the
Brauer--Manin analysis of Bright and Loughran
(\href{https://arxiv.org/abs/1908.02526}{arXiv:1908.02526}).  These works
provide strong structure, average abundance, and obstruction information, but
none supplies positive accepted mass for each remaining prime.

Consequently the unresolved statement is still pointwise:
\[
  \forall p\ \exists\text{ a sound SWC generator with proved mass }
  \delta(p)>0.
\]
A complete family cover, direct identity, nonconstructive count, or compatible
descent theorem could discharge it.  Finite verification, density-one
solvability, average representation counts, and improved conditional mass do
not discharge it by themselves.

